\part{II: Worlds and Admissibility}\label{part-ii-worlds-and-admissibility}

\hypertarget{chapter-5-possible-worlds-as-admissible-states}{%
\chapter*{Chapter 5 --- Possible Worlds as Admissible States}
\addcontentsline{toc}{chapter}{Chapter 5 --- Possible Worlds as Admissible States}
\markboth{Chapter 5 --- Possible Worlds as Admissible States}{Chapter 5 --- Possible Worlds as Admissible States}\label{chapter-5-possible-worlds-as-admissible-states}}

Take the most conservative starting point available: a \textbf{world} \(w\) is a
state satisfying some governing regime \(R\) (Chapter 1) ---
\(w \in \operatorname{Adm}(R)\). As a first approximation, modal necessity
is invariance across the admissible states \(R\) reaches:
\[\Box_R p \quad\text{iff}\quad p \text{ holds throughout the } R\text{-accessible admissible states}.\]

This formula is provisional, and it's worth being honest about why
before using it for anything: it makes a single relation, \(R\), do two
jobs at once --- deciding which states are \emph{admissible at all}, and
deciding which admissible states are \emph{relevant to evaluating \(\Box\) at a
given world}. Those are not obviously the same question. Whether a state
satisfies the constraints governing the domain (admissibility) and
whether that state is one this particular world can ``see'' for purposes
of modal evaluation (accessibility) are different facts about a state,
and a definition that quietly identifies them is buying convenience at
the cost of a real distinction. Chapter 6 exists specifically to undo
that conflation; this chapter's job is only to introduce
\(\operatorname{Adm}(R)\) itself; readers should treat the box formula
above as scaffolding to be dismantled, not as this chapter's finished
result.

\textbf{What admissibility does and doesn't settle.} \(\operatorname{Adm}(R)\)
gives worlds a structural role --- the space of states a regime permits ---
without committing to what worlds metaphysically \emph{are}. Priest's own
survey of realism, actualism, and Meinongian approaches to possible
worlds is a debate about the second question; this framework only needs
an answer to the first one to get started, and Priest introduces
possible-world semantics~\cite{priest2008} in exactly this same methodologically modest
role before taking up the ontological question separately. Nothing here
forces a choice between treating admissible states as concretely existing
alternatives, as actualized constructions, or as objects that need not
exist in the same sense actual things do. The chapter deliberately leaves
that choice open, because nothing built so far depends on it being made.

\hypertarget{chapter-6-necessity-accessibility-and-invariance}{%
\chapter*{Chapter 6 --- Necessity, Accessibility, and Invariance}
\addcontentsline{toc}{chapter}{Chapter 6 --- Necessity, Accessibility, and Invariance}
\markboth{Chapter 6 --- Necessity, Accessibility, and Invariance}{Chapter 6 --- Necessity, Accessibility, and Invariance}\label{chapter-6-necessity-accessibility-and-invariance}}

Chapter 5's box formula conflated two relations. This chapter separates
them and, in doing so, corrects something the book's own outline had
gotten slightly wrong: ``necessity as invariance,'' stated without further
qualification, describes something looser than what Kripke semantics~\cite{kripke1963}
actually is. Kripke necessity is not invariance across admissible states
generally --- it's truth throughout the states a specific accessibility
relation selects, and that relation is additional structure, not a
byproduct of admissibility itself.

\textbf{Two parameters.} Let \(R\) continue to determine admissibility --- which
states satisfy the governing regime --- and introduce \(\mathcal{R}\),
independently, as the modal accessibility relation: \(w\mathcal{R}v\) means
\(v\) is relevant to evaluating modal claims at \(w\). Necessity becomes
\[w \models \Box p \quad\iff\quad \forall v\, (w \mathcal{R} v \Rightarrow v \models p),\]
with no requirement that \(\mathcal{R}\) be derived from, or even
compatible with, \(R\) in any fixed way. A state can be \(\mathcal{R}\)-related
to \(w\) without being \(R\)-admissible in the same sense \(w\) is (worlds
under consideration for modal purposes need not all satisfy identical
constraints), and a state can be \(R\)-admissible without being
\(\mathcal{R}\)-accessible from any world in particular (an admissible
state simply not under consideration from a given vantage point).

\begin{figure}[htbp]
\centering
\begin{tikzpicture}[
  font=\small,
  adm/.style={draw, circle, minimum size=7mm, inner sep=0pt},
  inadm/.style={draw, circle, dashed, gray, minimum size=7mm, inner sep=0pt},
  acc/.style={-{Stealth[length=2.2mm]}}
]
\node[adm]   (w1) at (0,0)      {$w_1$};
\node[adm]   (w2) at (2.4,0.9)  {$w_2$};
\node[inadm] (w3) at (4.8,0)    {$w_3$};
\node[adm]   (w4) at (2.4,-1.5) {$w_4$};
\node[inadm] (w5) at (0,-2.6)   {$w_5$};
\node[adm]   (w6) at (4.8,-2.6) {$w_6$};
\draw[acc] (w1) -- (w2);
\draw[acc] (w2) -- (w3);
\draw[acc] (w1) -- (w4);
\draw[acc] (w4) -- (w6);
\draw[acc] (w5) -- (w4);
\draw[acc] (w3) to[bend left=25] (w6);
\end{tikzpicture}
\caption{Admissibility and accessibility are predicates of different
types. Arrows are accessibility edges $w \mathcal{R} v$; solid nodes are
$R$-admissible states, dashed gray nodes are not. $w_2 \mathcal{R} w_3$
holds even though $w_3 \notin \operatorname{Adm}(R)$ --- an accessible
world need not be admissible --- while $w_6 \in \operatorname{Adm}(R)$ is
reached only from an inadmissible state and from $w_4$; admissibility
generates no accessibility edge of its own. The rejected single-$R$
formulation of Chapter 5 would have made this frame inexpressible.}
\label{fig:adm-vs-acc}
\end{figure}

\textbf{What this buys.} Priest's various normal modal systems --- the
differences between K, T, S4, S5 and the rest --- are differences in what
structural properties \(\mathcal{R}\) is required to have: reflexivity,
transitivity, symmetry, and combinations of these, imposed on
accessibility specifically. None of this needs to touch \(R\) at all. A
system can strengthen or weaken its accessibility structure while
leaving admissibility exactly as it was, or vice versa --- change what
counts as an admissible state without touching which admissible states
are mutually accessible. Treating ``changing the modal logic'' and
``changing the admissibility regime'' as the same kind of move would have
been a real error, not a harmless simplification, and separating them
here is what makes the next chapter's argument possible at all.

\hypertarget{chapter-7-non-normal-and-impossible-worlds}{%
\chapter*{Chapter 7 --- Non-Normal and Impossible Worlds}
\addcontentsline{toc}{chapter}{Chapter 7 --- Non-Normal and Impossible Worlds}
\markboth{Chapter 7 --- Non-Normal and Impossible Worlds}{Chapter 7 --- Non-Normal and Impossible Worlds}\label{chapter-7-non-normal-and-impossible-worlds}}

This chapter is Part II's centerpiece, and it starts by retracting a
claim rather than building on one. An earlier working note in this
project's correspondence matrix rated ``a non-normal world deletes the
inference rule itself from what's admissible'' as a direct match to this
book's refusal machinery. That rating is downgraded here, from Direct to
Bridge, until this chapter actually establishes it --- asserting a
correspondence because the vocabulary sounds compatible is exactly the
kind of move this book has tried not to make, and there's no reason to
exempt this claim from that discipline.

\textbf{What Priest's non-normal worlds actually do, reconstructed first.} At
an ordinary (``normal'') world, \(\Box p\) is evaluated compositionally: true
iff \(p\) holds at every \(\mathcal{R}\)-accessible world, per Chapter 6. At
a \textbf{non-normal} world, that compositional rule is suspended --- the truth
value of \(\Box p\) at a non-normal world is not computed from what holds
at accessible worlds at all; it can be assigned independently, including
being false at a non-normal world even for a \(p\) that is a logical truth
holding at every other world. The consequence that matters most: modus
ponens itself, ordinarily a validity-preserving step everywhere, need
not hold at a non-normal world, because the world's evaluation behavior
is no longer tied to the mechanism that made modus ponens valid
elsewhere. \textbf{Impossible worlds} generalize this further still: worlds at
which the ordinary logical laws governing the rest of the framework may
simply fail to apply, not merely modal ones.

\textbf{What actually reproduces this, once reconstructed.} The behavior
above is not a targeted decision to decline one specific transformation
--- that would be refusal, in Chapter 4's sense, and refusal is a
\emph{Level 3} operation: a constitutive choice about which repair moves an
agent or system will make. What non-normal worlds do is different in
kind: the \emph{evaluation rule itself} is locally not in force. In this
book's vocabulary, that is best captured not by refusal but by treating
the regime as state-indexed rather than global: \(R_w\), a regime that can
itself vary from state to state, rather than one fixed \(R\) applied
everywhere with exceptions carved out afterward. A non-normal world is a
state whose local \(R_w\) simply does not include the closure condition
(Chapter 3) that would make \(\Box\)'s compositional rule, or modus ponens,
admissible moves there. This is a real correspondence, but a different
and more specific one than the original matrix entry claimed --- and it's
one this chapter has now actually earned, rather than assumed.

\textbf{The control case Part III supplies.} Before Part III existed, it
would have been tempting to treat ``impossible'' and ``contains a
contradiction'' as roughly the same idea. Part III makes that identification
untenable. An FDE world can contain both \(p\) and \(\neg p\) --- a glut, in
Chapter 10's sense --- while its consequence behavior remains perfectly
well-defined: Chapter 11's restricted propagation tells you exactly what
follows from that glut and what doesn't, compositionally, with no
suspension of any evaluation rule anywhere. Such a world is inconsistent
without being impossible in the sense this chapter has just defined, and
without being non-normal --- nothing about its evaluation machinery has
been switched off. Conversely, a non-normal or impossible state can fail
to validate modus ponens, or fail to respect \(\Box\)'s compositional rule,
while containing no glut at all --- every atomic proposition perfectly
determinate, no overdetermination anywhere, and still a state where some
structural rule simply isn't in force. These are three independent axes,
not one phenomenon under three names:

\[\boxed{\text{inconsistency} \neq \text{impossibility} \neq \text{non-normality}}\]

This belongs beside Chapter 3's four-way consistency/closure/coherence
split as the same kind of discipline applied to modal vocabulary instead
of evidential vocabulary. And it produces one sentence worth treating as
foundational for the rest of this book's modal material:

\begin{quote}
\textbf{Impossibility is indexed to a constraint regime; contradiction is
indexed to support.}
\end{quote}

A state is impossible or not relative to which \(R_w\) it's being measured
against --- the same state can be admissible under one regime and
inadmissible under a stricter one, with nothing about its content having
changed. A proposition is contradictory or not as a fact about its
support, full stop, independent of any regime. Confusing these two
coordinates --- treating the presence of a glut as automatically making a
world impossible --- would undo exactly the distinction Part III spent
four chapters establishing.

\hypertarget{chapter-8-conditionals-as-constrained-transport}{%
\chapter*{Chapter 8 --- Conditionals as Constrained Transport}
\addcontentsline{toc}{chapter}{Chapter 8 --- Conditionals as Constrained Transport}
\markboth{Chapter 8 --- Conditionals as Constrained Transport}{Chapter 8 --- Conditionals as Constrained Transport}\label{chapter-8-conditionals-as-constrained-transport}}

This chapter combines what earlier planning had kept as two separate
chapters --- strict implication and Lewis-style conditional logic~\cite{lewis1973} --- because
both are attempts at the same underlying question, and treating them
separately obscured that.

\textbf{The question.} Given a state \(w\) and a supposition \(p\) --- not
necessarily true at \(w\) --- which transformations away from \(w\) are
licensed when evaluating a further claim \(q\) under that supposition?
Strict implication answers this with maximum severity: \(p\) strictly
implies \(q\) iff \(q\) holds at every state where \(p\) holds, full stop,
which is exactly the source of strict implication's well-known paradoxes
(a necessary truth is strictly implied by anything, an impossible
proposition strictly implies anything) --- the same over-permissive
propagation problem Chapter 3 and Chapter 11 have already diagnosed in
other settings, showing up again here in conditional form. Lewis and
Stalnaker's similarity-sphere semantics answers it with an ordering:
evaluate \(q\) at the \(p\)-worlds \emph{most similar} to \(w\), rather than at all
of them.

\textbf{A candidate reformulation.} In this book's vocabulary, the question
becomes: given \(w\) and \(p\), which \(p\)-satisfying transformations of \(w\)
are \(R\)-admissible? Write this as
\[p \Rightarrow_R q\]
with truth depending on admissible \(p\)-transformations of \(w\) rather than
on material implication or on an externally supplied similarity metric.
This reformulation has real content --- it ties the conditional's behavior
to the same admissibility apparatus governing everything else in this
book, rather than treating conditionals as needing their own separate
primitive machinery --- but it should not be oversold. What it does not
yet do is tell you \emph{which} transformations count as admissible for a
given supposition, which is exactly the work Lewis's similarity ordering
was doing, and which this book's distinguishability geometry has been
suggested, elsewhere, as a candidate replacement for.

\textbf{The debt this chapter leaves open, deliberately.} It would be
premature to claim distinguishability geometry \emph{replaces} Lewisian
similarity spheres --- the correspondence matrix marks that Bridge, and
this chapter doesn't upgrade it, because nothing here has shown the
similarity ordering can actually be derived rather than separately
stipulated. The honest question to end on: can a similarity ordering
over \(p\)-satisfying transformations be derived from the geometry of
admissible transformation itself --- from the same structure that already
determines \(\operatorname{Adm}(R)\) and \(R_w\) --- rather than supplied as
independent primitive structure the way Lewis's own semantics requires?
If yes, conditionals stop needing a bespoke ordering relation and become
a further consequence of admissibility geometry already on hand. If no,
conditionals are a genuinely independent piece of structure this
framework hasn't reduced, and the book should say so plainly rather than
paper over it. This chapter poses that question; it does not answer it.

\hypertarget{status-of-this-draft}{%
\section{Status of this draft}\label{status-of-this-draft}}

Part II's progression --- alternative states, accessibility, violations of
admissibility, transport under supposition --- holds together, and the
chapter that mattered most going in, Chapter 7, survived contact with
Part III's control case rather than needing to be softened to fit it.
One real result came out of testing rather than assuming: the original
correspondence-matrix claim about non-normal worlds and refusal was
wrong in its specifics, even though something in that neighborhood is
true --- the corrected version (state-indexed regime \(R_w\), not refusal)
is more precise and, on inspection, more defensible than the claim it
replaced. Chapter 8 ends with a named open debt rather than a premature
claim, per instruction. Worth checking once Part V is drafted: Chapter 7's
\(R_w\) move (regimes varying by state) is very close to what Part V's
identity-under-transport material will need for tracking a single
entity across states with potentially different local constraints ---
that connection wasn't planned in advance and should be verified rather
than assumed to hold.
